\clearpage
\section*{What substitute ordinances change}
\textit{September 30, 2026. Agent-drafted work on branch \texttt{explore-alderman-measures}; Jacob asked for a measure
of how often and how much aldermen cut projects before approving them, as Hopkins was said to have done.}

Nearly half of the applications that passed in eLMS (2,237 of about 4,700) have both the ordinance introduced and a
substitute that replaced it, and eLMS keeps both. \path{tasks/compare_substitute_ordinances} compares what the two
state: the district after (by its highest floor-area ratio) and, from each version's Type 1 narrative or planned
development's bulk table, the project's floor-area ratio, dwelling units and height. Both versions are read by the
same rules and a pair is compared only where they are read alike: the same rule in both, the Clerk's text layer and
tesseract agreeing, and no tenfold gap or implausible ratio, since in a sample the first rules' apparent changes were
often misreadings, which crowd into the few pairs that differ. Against hand reads of 30 applications drawn at random
(\path{tasks/audits/zoning_record_validation}, section 10), 25 of the 26 changes the comparison finds are real and 26
of its 28 unchanged pairs are unchanged; the rules were corrected once after the first scoring, which confirmed 26 of
29 changes.

1,145 passed applications are compared on at least one field. Most substitutes leave the project as introduced, and
those that change it do not change it in one direction: 94 of the 1,137 with an alderman were made smaller, 90 larger
and 7 both, and 946 kept the same district, ratio, units and height. Planned developments change more often (31
percent, against 15 percent). Aldermen do not differ in this beyond chance: across the 31 with at least 10 compared
applications, the shares made smaller, made larger and changed vary as much as reassigning applications at random
within kind and period does (permutation p = 0.34, 0.12 and 0.33; \path{tasks/audits/zoning_stalls_and_delays}).
Hopkins's substitutes made 5 of his 44 compared applications smaller and 1 larger. That is consistent with negotiation
being settled before an application is filed, or showing in what is filed at all, rather than in the substitute. A
project an alderman blocks outright, as Old Town Canvas was, never passes and is not in this comparison.

\textit{Open.} Half of the passed applications with both versions are not compared, mostly because the substitute is
the ordinance alone; only the first three pages of the narratives attached from 2023 were read, so recent planned
developments' bulk tables are mostly missed (Conway has 7 compared), and the share made smaller falls to 2 percent of
applications from 2023, when the versions come from those attachments, which may be a change of source rather than of
practice. Reading the attachments in full would extend the comparison.
